%% Template for the submission to:
%%   The Annals of Statistics [AOS]
%%
%%%%%%%%%%%%%%%%%%%%%%%%%%%%%%%%%%%%%%%%%%%%%%
%% In this template, the places where you   %%
%% need to fill in your information are     %%
%% indicated by '???'.                      %%
%%                                          %%
%% Please do not use \input{...} to include %%
%% other tex files. Submit your LaTeX       %%
%% manuscript as one .tex document.         %%
%%%%%%%%%%%%%%%%%%%%%%%%%%%%%%%%%%%%%%%%%%%%%%

\documentclass{imsart}

%% Packages
\RequirePackage{amsthm,amsmath,amsfonts,amssymb,booktabs,array}
\RequirePackage[numbers]{natbib}
%\RequirePackage[authoryear]{natbib}%% uncomment this for author-year citations
\RequirePackage[colorlinks,citecolor=blue,urlcolor=blue]{hyperref}%% uncomment this for coloring bibliography citations and linked URLs
\RequirePackage{graphicx}%% uncomment this for including figures

\startlocaldefs
%%%%%%%%%%%%%%%%%%%%%%%%%%%%%%%%%%%%%%%%%%%%%%
%%                                          %%
%% Uncomment next line to change            %%
%% the type of equation numbering           %%
%%                                          %%
%%%%%%%%%%%%%%%%%%%%%%%%%%%%%%%%%%%%%%%%%%%%%%
%\numberwithin{equation}{section}
%%%%%%%%%%%%%%%%%%%%%%%%%%%%%%%%%%%%%%%%%%%%%%
%%                                          %%
%% For Axiom, Claim, Corollary, Hypothesis, %%
%% Lemma, Theorem, Proposition              %%
%% use \theoremstyle{plain}                 %%
%%                                          %%
%%%%%%%%%%%%%%%%%%%%%%%%%%%%%%%%%%%%%%%%%%%%%%
\theoremstyle{plain}
\newtheorem{thm}{\textsc{Theorem}}[section]
\newtheorem{prop}[thm]{\textsc{Proposition}}
\newtheorem{cor}[thm]{\textsc{Corollary}}
\newtheorem{lem}[thm]{\textsc{Lemma}}
%%%%%%%%%%%%%%%%%%%%%%%%%%%%%%%%%%%%%%%%%%%%%%
%%                                          %%
%% For Assumption, Definition, Example,     %%
%% Notation, Property, Remark, Fact         %%
%% use \theoremstyle{remark}                %%
%%                                          %%
%%%%%%%%%%%%%%%%%%%%%%%%%%%%%%%%%%%%%%%%%%%%%%
%\theoremstyle{remark}
%\newtheorem{???}{???}
%\newtheorem*{???}{???}
%\newtheorem{???}{???}[???]
%\newtheorem{???}[???]{???}
%%%%%%%%%%%%%%%%%%%%%%%%%%%%%%%%%%%%%%%%%%%%%%
%% Please put your definitions here:        %%
%%%%%%%%%%%%%%%%%%%%%%%%%%%%%%%%%%%%%%%%%%%%%%
\newcommand{\david}[1]{{\small{\color{blue} \sf $\heartsuit$ [David: #1]}}}
\newcommand{\piotr}[1]{{\small{\color{red} \sf $\clubsuit$ Piotr: #1}}}
\newcommand{\cpz}[1]{{\color{red} #1}}
\newcommand{\paul}[1]{{\small{\color{green} \sf $\spadesuit$ Paul: #1}}}
\newcommand{\E}{\mathbb{E}}
\newcommand{\N}{\mathbb{N}}
\newcommand{\R}{\mathbb{R}}
\newcommand{\I}{\mathbb{I}}
\newcommand{\M}{\mathcal{M}}
\newcommand{\T}{\mathcal{T}}
\newcommand{\bs}{\boldsymbol}
\newcolumntype{M}[1]{>{\centering\arraybackslash}m{#1}}
\endlocaldefs

\begin{document}

\begin{frontmatter}
%%%%%%%%%%%%%%%%%%%%%%%%%%%%%%%%%%%%%%%%%%%%%%
%%                                          %%
%% Enter the title of your article here     %%
%%                                          %%
%%%%%%%%%%%%%%%%%%%%%%%%%%%%%%%%%%%%%%%%%%%%%%
\title{Cover letter to "Improving variable selection properties by leveraging external data"}
%\title{A sample article title with some additional note\thanksref{T1}}
%\runtitle{Improving variable selection properties by using external data}
%\thankstext{T1}{A sample of additional note to the title.}

\begin{aug}
%%%%%%%%%%%%%%%%%%%%%%%%%%%%%%%%%%%%%%%%%%%%%%%
%% Only one address is permitted per author. %%
%% Only division, organization and e-mail is %%
%% included in the address.                  %%
%% Additional information can be included in %%
%% the Acknowledgments section if necessary. %%
%% ORCID can be inserted by command:         %%
%% \orcid{0000-0000-0000-0000}               %%
%%%%%%%%%%%%%%%%%%%%%%%%%%%%%%%%%%%%%%%%%%%%%%%
% \author[A]{\fnms{Paul}~\snm{Rognon-Vael}\ead[label=e1]{paul.rognon@gmail.com}},
% \author[A]{\fnms{David}~\snm{Rossell}\ead[label=e2]{rosselldavid@gmail.com}}
% \and
% \author[B]{\fnms{Piotr}~\snm{Zwiernik}\ead[label=e3]{piotr.zwiernik@utoronto.ca}}
% %%%%%%%%%%%%%%%%%%%%%%%%%%%%%%%%%%%%%%%%%%%%%%
% %% Addresses                                %%
% %%%%%%%%%%%%%%%%%%%%%%%%%%%%%%%%%%%%%%%%%%%%%%
% \address[A]{Department of Economics and Business, Universitat Pompeu Fabra \printead[presep={ ,\ }]{e1,e2}}
% \address[B]{Department of Statistical Sciences, University or Toronto
% \printead[presep={,\ }]{e3}}
\end{aug}





% \begin{keyword}[class=MSC]
% \kwd[Primary ]{62F07}
% \kwd[; secondary ]{62C12}
% \kwd{62R07}
% \end{keyword}

% \begin{keyword}
% \kwd{variable selection}
% \kwd{high-dimensional statistics}
% \kwd{$\ell_0$ penalty}
% \kwd{empirical Bayes}
% \end{keyword}
\end{frontmatter}
%%%%%%%%%%%%%%%%%%%%%%%%%%%%%%%%%%%%%%%%%%%%%%
%% Please use \tableofcontents for articles %%
%% with 50 pages and more                   %%
%%%%%%%%%%%%%%%%%%%%%%%%%%%%%%%%%%%%%%%%%%%%%%
%\tableofcontents

%%%%%%%%%%%%%%%%%%%%%%%%%%%%%%%%%%%%%%%%%%%%%%
%%%% Main text entry area:

%Please find attached our paper "Improving variable selection properties by using external data". 
Our work studies how one can push the mathematical limits for model recovery (structural learning), when one has access to external information. 
Briefly, there is a growing methodological and applied literature that uses data integration (of various types) to improve inference, particularly in settings where one has scarce data. For example, when studying a rare cancer one may leverage on related cancer (sub)types to better identify biomarkers. As another example, transfer learning seeks to extrapolate findings from a (typically large) training data to new settings where data is more scarce. There are many other examples, briefly reviewed in the paper.
These strategies are popular and have shown substantial empirical gains.

Our main contribution is providing a theoretical backbone for these gains, in the setting of sparse support recovery.
%Briefly, most high-dimensional theory assumes that one only has access to the data that's of immediate interest, for example a dataset with the response and covariate values for $n$ individuals. In practice one also has external information, for example in the form of meta-covariates such as whether each covariate was found likely to have an effect in a related previous study (from a different cancer type, say). There are many such examples and the applied literature showed that, by integrating such external information into the data analysis, one may obtain significant empirical gains. 
We study what gains are attainable %why and how these gains occurs 
in a simplified setting where the external data that is being integrated allows splitting %meta-covariates split 
the parameters into blocks. We consider the Gaussian sequence model and high-dimensional linear regression. First we consider an oracle point-of-view that portrays the best possible gains one could wish for. Second, we propose an empirical Bayes data analysis method. When the external data is sufficiently informative it is possible to achieve consistent model selection where otherwise it's provably impossible, and even when consistency would be possible one may obtain better rates. To our knowledge our work is the first studying how external information may aid model selection (aka structural learning) consistency, complementing previous studies related to hypothesis testing or information-theoretic gains (as discussed in the introduction).\\

Thank you for considering our work and best regards,\\

\noindent Paul Rognon-Vael, David Rossell \& Piotr Zwiernik
\end{document}
